% Source PDF: source/普通高考/2015/2015安徽文.pdf, page 1, question 7.
% pdfimages -list found no embedded bitmap; the program flowchart is vector line art.
% Grid evidence: tmp/redraw_flowchart_2015_anhui_liberal/grid/page1_grid100.png and q7_block_grid50.png.
% Comparison crop: tmp/redraw_flowchart_2015_anhui_liberal/crops/q7_orig_final.png,
% taken from the ungridded 300 dpi page render with a safety border.
% Elements: rounded start/end terminals; a=1,n=1 initialization rectangle; the
% |a-1.414|>=0.005? decision diamond; update rectangle a=1+1/(1+a); n=n+1;
% output parallelogram; branch labels 是/否; and the left loop-back connector.
% Relations: start -> a=1,n=1 -> decision; 是 -> update -> n=n+1 -> decision;
% 否 -> 输出 n -> 结束. The initialization stem and loop return meet before the
% decision; the loop's horizontal segment points right, followed by one downward
% arrow into the diamond. No dashed segments or shared output-loop connections.
\begin{tikzpicture}[
  x=1cm,y=1cm,baseline,
  flowarrow/.style={-{Latex[length=2.1mm,width=1.5mm]},line width=0.45pt},
  flowtext/.style={anchor=center,align=center,
    execute at begin node={\setlength{\parindent}{0pt}}},
  every node/.style={flowtext},
  nodebox/.style={flowtext,draw,line width=0.45pt,inner sep=2pt,minimum height=0.66cm},
  branchlabel/.style={flowtext,inner sep=0pt},
  term/.style={nodebox,rounded corners=2.5pt,minimum width=1.35cm,text width=1.35cm},
  io/.style={nodebox,shape=trapezium,trapezium left angle=70,trapezium right angle=110,
    minimum width=1.90cm,text width=1.60cm,minimum height=0.72cm},
  decision/.style={nodebox,diamond,aspect=4.0,minimum width=4.95cm,minimum height=1.15cm,inner sep=1pt}
]
  \node[term] (start) at (0,9.30) {开始};
  \node[nodebox,text width=3.20cm,minimum width=3.20cm] (init) at (0,8.20) {$a=1,n=1$};
  \node[decision] (test) at (0,6.25) {$|a-1.414|\geq 0.005?$};
  \node[nodebox,text width=3.55cm,minimum width=3.55cm,minimum height=1.05cm] (update) at (0,4.35)
    {$a=1+\dfrac{1}{1+a}$};
  \node[nodebox,text width=2.55cm,minimum width=2.55cm] (increment) at (0,2.70) {$n=n+1$};
  \node[io] (output) at (4.10,4.85) {输出 $n$};
  \node[term] (finish) at (4.10,3.20) {结束};

  % Main path and affirmative branch.
  \draw[flowarrow] (start.south) -- (init.north);
  % The initial stem joins the loop junction without an extra arrowhead.
  % The source keeps the horizontal return arrowhead clearly above the
  % separate downward arrowhead entering the decision.
  \draw (init.south) -- (0,7.60);
  \draw[flowarrow] (0,7.60) -- (test.north);
  \draw[flowarrow] (test.south) -- node[right=2pt,inner sep=0pt] {是} (update.north);
  \draw[flowarrow] (update.south) -- (increment.north);

  % Negative branch to output and termination, independent of the loop route.
  \draw[flowarrow] (test.east) -- node[above=2pt,inner sep=0pt] {否}
    (4.10,6.25) -- (output.north);
  \draw[flowarrow] (output.south) -- (finish.north);

  % Return path from n=n+1 to the single input point above the decision.
  \draw (increment.south) -- (-3.55,2.30) -- (-3.55,7.60);
  \draw[flowarrow] (-3.55,7.60) -- (0,7.60);

\end{tikzpicture}
